% Preámbulo común del material docente · XeLaTeX
\documentclass[11pt,a4paper]{article}
\usepackage[margin=22mm,top=20mm,bottom=22mm]{geometry}
\usepackage{fontspec}
\setmainfont{Avenir Next}[BoldFont={Avenir Next Demi Bold}, BoldItalicFont={Avenir Next Demi Bold Italic}]
\setsansfont{Avenir Next}
\setmonofont{Menlo}[Scale=0.88]
\newfontfamily\symfont{Apple Symbols}
\newfontfamily\unifont{Arial Unicode MS}
\usepackage{unicode-math}
\setmathfont{latinmodern-math.otf}
\usepackage{polyglossia}\setdefaultlanguage{spanish}
\usepackage{xcolor,booktabs,tabularx,enumitem,parskip,microtype,titlesec,fancyhdr}
\usepackage[most]{tcolorbox}
\usepackage[hidelinks]{hyperref}
\emergencystretch=2em

\definecolor{tinta}{HTML}{1F2933}
\definecolor{acento}{HTML}{B7410E}
\definecolor{suave}{HTML}{F4F1EC}
\definecolor{gris}{HTML}{6B7280}
\definecolor{linea}{HTML}{D6D0C6}
\color{tinta}

\titleformat{\section}{\Large\bfseries\color{acento}}{\thesection}{0.6em}{}
\titleformat{\subsection}{\large\bfseries}{\thesubsection}{0.6em}{}
\titlespacing*{\section}{0pt}{1.6em}{0.5em}
\titlespacing*{\subsection}{0pt}{1.1em}{0.3em}
\setlist{nosep, leftmargin=1.4em, itemsep=2pt}
\renewcommand{\arraystretch}{1.25}

% bloque de órdenes de terminal
\newtcblisting{orden}{listing only, colback=suave, colframe=linea, boxrule=0.4pt, arc=2pt,
  left=6pt, right=6pt, top=4pt, bottom=4pt, listing options={basicstyle=\ttfamily\small, breaklines=true, columns=fullflexible, keepspaces=true}}
% nota destacada
\newtcolorbox{nota}[1][Nota]{colback=white, colframe=acento, boxrule=0.6pt, arc=2pt, title=#1,
  fonttitle=\bfseries\small, coltitle=white, colbacktitle=acento, left=6pt, right=6pt}
% preguntas
\newcommand{\pregunta}[2]{\item[\textbf{(#1)}] #2}
\newenvironment{preguntas}{\begin{description}[leftmargin=2.4em, style=nextline, itemsep=4pt]}{\end{description}}
\newcommand{\ok}{{\symfont ✔}}
\newcommand{\nok}{{\symfont ✘}}
\newcommand{\qed}{{\symfont ∎}}
\newcommand{\codigo}[1]{\texttt{#1}}

\pagestyle{fancy}\fancyhf{}
\renewcommand{\headrulewidth}{0pt}
\fancyfoot[C]{\small\color{gris}\docpie\ · página \thepage}

\newcommand{\cabecera}[3]{%
  {\Huge\bfseries\color{acento}#1\par}\vspace{4pt}
  {\large #2\par}\vspace{2pt}
  {\color{gris}#3\par}\vspace{8pt}
  {\color{linea}\hrule height 0.6pt}\vspace{10pt}}

\newcommand{\docpie}{La terminal desde cero}
\begin{document}
\cabecera{La terminal desde cero}{Métodos numéricos · guía para la primera sesión}{Léela antes de tocar nada. Diez minutos.}

\section{Qué es la terminal}

Todo lo que haces con el ratón (abrir una carpeta, abrir un programa, crear un archivo) se puede hacer \textbf{escribiendo órdenes}. El programa donde se escriben esas órdenes es la \textbf{terminal}: una ventana negra (o azul) con un cursor parpadeando. Parece de hacker; no lo es. Es solo otra forma de decirle al computador qué hacer, más precisa y más rápida.

Una orden se escribe, se pulsa \textbf{Enter}, y el computador responde con texto. Eso es todo.

¿Por qué la usamos en este curso? Porque los programas que vamos a correr no tienen botones. Se ejecutan con una orden, escriben su resultado en la terminal y guardan archivos en una carpeta. Así trabaja casi toda la matemática computacional.

\section{Abrir la terminal}

\textbf{Windows (los PCs del laboratorio):} pulsa la tecla Windows, escribe \codigo{terminal} y abre \textbf{Terminal} (o \textbf{Windows PowerShell} si no aparece). Maximiza la ventana: la vamos a necesitar grande.

\textbf{Mac:} Cmd + espacio, escribe \codigo{terminal}, Enter.

Verás algo así:
\begin{orden}
PS C:\Users\estudiante>
\end{orden}
Eso se llama el \textbf{prompt}. Te dice \textbf{en qué carpeta estás} (\codigo{C:\textbackslash Users\textbackslash estudiante}) y espera una orden. En Mac se ve parecido pero con \codigo{/Users/nombre} y termina en \codigo{\%} o \codigo{\$}.

\section{Las cinco órdenes que necesitas}

Escribe cada una y pulsa Enter. Fíjate en lo que responde.

\begin{tabularx}{\linewidth}{@{}l X l@{}}
\toprule
\textbf{orden} & \textbf{qué hace} & \textbf{ejemplo} \\
\midrule
\codigo{pwd} & te dice en qué carpeta estás (\emph{print working directory}) & \codigo{pwd} \\
\codigo{ls} & lista lo que hay en la carpeta actual (\emph{list}) & \codigo{ls} \\
\codigo{cd nombre} & entra en la carpeta \codigo{nombre} (\emph{change directory}) & \codigo{cd Documents} \\
\codigo{cd ..} & sube a la carpeta de arriba & \codigo{cd ..} \\
\codigo{mkdir nombre} & crea una carpeta & \codigo{mkdir metodos} \\
\bottomrule
\end{tabularx}

Tres trucos que ahorran mucho sufrimiento:
\begin{itemize}
\item \textbf{Tab} completa nombres. Escribe \codigo{cd Doc} y pulsa Tab: se convierte en \codigo{cd Documents}.
\item \textbf{Flecha arriba} repite la orden anterior. No vuelvas a escribir nada largo.
\item Si escribes algo mal, la terminal te lo dice en rojo o con «no se reconoce». No pasa nada. Corrige y vuelve a intentar. \textbf{No puedes romper el computador con estas órdenes.}
\end{itemize}

Prueba ahora:
\begin{orden}
cd Documents
mkdir metodos
cd metodos
pwd
\end{orden}
Deberías ver \codigo{C:\textbackslash Users\textbackslash <tu usuario>\textbackslash Documents\textbackslash metodos}. Esa es tu carpeta de trabajo del curso.

\section{Qué significa «correr código»}

Un programa es un archivo de texto con instrucciones. \textbf{Python} es el programa que lee esas instrucciones y las ejecuta. Cuando en la terminal escribes
\begin{orden}
python archivo.py
\end{orden}
le estás diciendo: «Python, lee \codigo{archivo.py} y haz lo que dice». Python trabaja, escribe resultados en la terminal, y cuando termina vuelve a aparecer el prompt.

Comprueba que Python está instalado:
\begin{orden}
python --version
\end{orden}
Debe responder \codigo{Python 3.1x.x}. Si dice que no se reconoce, avisa al docente (en Windows a veces la orden es \codigo{py} en vez de \codigo{python}; prueba \codigo{py --version}).

\section{Descargar el material del curso}

El código está en GitHub, una página donde se comparten programas. No necesitas cuenta.
\begin{enumerate}
\item Abre en el navegador: \url{https://github.com/cmorregof/rl_metodos}
\item Botón verde \textbf{Code} $\to$ \textbf{Download ZIP}.
\item El archivo \codigo{rl\_metodos-main.zip} queda en \codigo{Descargas} (\codigo{Downloads}). Descomprímelo (clic derecho $\to$ \emph{Extraer todo}) y mueve la carpeta \codigo{rl\_metodos-main} a \codigo{Documents\textbackslash metodos}.
\end{enumerate}

Ahora en la terminal:
\begin{orden}
cd Documents\metodos\rl_metodos-main
ls
\end{orden}
Deberías ver las carpetas \codigo{01\_biseccion}, \codigo{02\_punto\_fijo}, … Si \codigo{ls} no las muestra, no estás en la carpeta correcta: usa \codigo{pwd} para ver dónde estás y \codigo{cd} para moverte.

(En Mac las barras van al revés: \codigo{cd Documents/metodos/rl\_metodos-main}.)

\section{Instalar el proyecto de hoy}

Cada carpeta es un proyecto independiente. Entra en el primero y crea un «entorno»: una carpeta privada donde Python instala lo que ese proyecto necesita sin tocar nada más del computador.
\begin{orden}
cd 01_biseccion
python -m venv .venv
.venv\Scripts\python -m pip install -e ".[dev]"
\end{orden}
(Mac: \codigo{.venv/bin/python -m pip install -e ".[dev]"}.)

La segunda orden tarda unos segundos y no dice nada. La tercera escribe muchas líneas; lo importante es que la última diga \codigo{Successfully installed ...}. Si aparece un error rojo, levanta la mano.

\begin{nota}[Desde ahora]
«El Python del proyecto» es \codigo{.venv\textbackslash Scripts\textbackslash python}. Cada orden del laboratorio empieza así. Nunca hace falta escribir \codigo{activate}.
\end{nota}

\section{Correr el agente}
\begin{orden}
.venv\Scripts\python -m bisectrl --seed 1
\end{orden}
Se abre una pantalla llena de paneles que se mueven. Es un programa aprendiendo en vivo. Dura uno o dos minutos. Si la ventana es pequeña se verá recortado: maximízala y vuelve a correr (flecha arriba, Enter).

Cuando termina, imprime un informe y dice dónde lo guardó: en la carpeta \codigo{runs}, en una subcarpeta con la fecha y hora. Ábrela desde el explorador de archivos y abre \codigo{informe.md} con el Bloc de notas (o con VS Code si está instalado). Es texto plano: se lee igual.

\section{Si algo se atasca}
\begin{itemize}
\item El programa no termina y quieres pararlo: \textbf{Ctrl + C}.
\item La terminal «se quedó rara»: ciérrala, abre otra y vuelve a la carpeta con \codigo{cd}.
\item «No se reconoce el comando»: revisa mayúsculas y que estés en la carpeta correcta (\codigo{pwd}).
\item Windows PowerShell dice algo de «ejecución de scripts deshabilitada»: es que escribiste \codigo{activate}; no hace falta. Usa siempre \codigo{.venv\textbackslash Scripts\textbackslash python}.
\end{itemize}

\section{Lo que acabas de aprender}

Sabes abrir una terminal, moverte por carpetas, comprobar que Python existe, crear un entorno para un proyecto y ejecutar un programa que escribe archivos. Eso es el 90\,\% de lo que hace un matemático computacional al empezar cualquier cosa. El otro 10\,\% es leer lo que el programa dice cuando falla, y eso se aprende fallando.
\end{document}
